\documentclass[11pt]{article}
\usepackage[utf8]{inputenc}
\usepackage{amsmath,amssymb,booktabs,geometry,verbatim}
\usepackage{multirow}
\usepackage{CJKutf8}
\geometry{margin=1in}
\usepackage[hidelinks]{hyperref}

\newcommand{\cn}[1]{\begin{CJK}{UTF8}{gbsn}#1\end{CJK}}

\title{CSP-S 2025 (Senior)}
\date{2025-10-25, 14:30--18:30 (4 hours)}

\begin{document}
\maketitle

\noindent
This file is tailored to QOJ: every problem reads from standard input and
writes to standard output (the original paper used the files
\texttt{<code>.in} / \texttt{<code>.out} named after each problem).

\noindent
2025 CCF Certified Software Professional, CSP-J/S 2025, second round,
Senior level (CSP-S).
Four problems, in contest order.
Time and memory limits are those shown on the QOJ problem pages.

\clearpage
\section*{Problem 1: Club Recruitment (\cn{社团招新})
  \hfill{\normalsize\href{https://qoj.ac/problem/14597}{QOJ 14597}}}

\noindent
Time limit: 1 s. \quad Memory limit: 512 MB. \quad Total points: 100.

\subsection*{Problem description}

Little L is a member of the school's algorithm society.
In this year's school club recruitment, Little L recruited $n$ new members
in total, where $n$ is \textbf{even}.
Now Little L wants to assign them to different departments of the society.

The algorithm society has three departments.
The satisfaction of the $i$-th ($1 \le i \le n$) new member with the
$j$-th ($1 \le j \le 3$) department is $a_{i,j}$.
The satisfaction of an assignment is defined as the sum, over all new
members, of their satisfaction with the department they are assigned to.
That is, if the $i$-th ($1 \le i \le n$) new member is assigned to the
$d_i$-th department, where $d_i \in \{1,2,3\}$, then the satisfaction of
the assignment is $\sum_{i=1}^{n} a_{i,d_i}$.

Little L does not want any single department to have too many new members.
Specifically, he requires that in the assignment \textbf{there is no}
department that is assigned \textbf{more than} $\frac{n}{2}$ new members.
You need to help Little L find the maximum satisfaction of an assignment
that satisfies his requirement.

\subsection*{Input}

\textbf{This problem contains multiple test cases.}

The first line of the input contains a positive integer $t$, the number of
test cases.

Then the test cases follow one after another.
For each test case:
\begin{itemize}
\item The first line contains a positive integer $n$, the number of new
  members.
\item Line $i+1$ ($1 \le i \le n$) contains three non-negative integers
  $a_{i,1}, a_{i,2}, a_{i,3}$, the satisfaction of the $i$-th new member
  with departments $1, 2, 3$ respectively.
\end{itemize}

\subsection*{Output}

For each test case, output one line containing one non-negative integer,
the maximum satisfaction of an assignment that satisfies Little L's
requirement.

\subsection*{Sample 1 input}

\begin{verbatim}
3
4
4 2 1
3 2 4
5 3 4
3 5 1
4
0 1 0
0 1 0
0 2 0
0 2 0
2
10 9 8
4 0 0
\end{verbatim}

\subsection*{Sample 1 output}

\begin{verbatim}
18
4
13
\end{verbatim}

\subsection*{Explanation of sample 1}

This sample contains three test cases.

For the first test case, the four new members can be assigned to
departments $1, 3, 1, 2$ respectively.
Then the numbers of new members in the three departments are $2, 1, 1$,
none exceeding $\frac{4}{2} = 2$, and the satisfaction is
$4 + 4 + 5 + 5 = 18$.

For the second test case, the four new members can be assigned to
departments $1, 1, 2, 2$ respectively.
Then the numbers of new members in the three departments are $2, 2, 0$,
none exceeding $\frac{4}{2} = 2$, and the satisfaction is
$0 + 0 + 2 + 2 = 4$.

For the third test case, the two new members can be assigned to
departments $2, 1$ respectively.
Then the numbers of new members in the three departments are $1, 1, 0$,
none exceeding $\frac{2}{2} = 1$, and the satisfaction is $9 + 4 = 13$.

\subsection*{Sample 2}

Large sample, not printed here;
provided as the downloadable files \texttt{club/club2.in} and \texttt{club/club2.ans}.

This sample satisfies the constraints of test cases $3, 4$.

\subsection*{Sample 3}

Large sample, not printed here;
provided as the downloadable files \texttt{club/club3.in} and \texttt{club/club3.ans}.

This sample satisfies the constraints of test cases $5 \sim 8$.

\subsection*{Sample 4}

Large sample, not printed here;
provided as the downloadable files \texttt{club/club4.in} and \texttt{club/club4.ans}.

This sample satisfies the constraints of test case $9$.

\subsection*{Sample 5}

Large sample, not printed here;
provided as the downloadable files \texttt{club/club5.in} and \texttt{club/club5.ans}.

This sample satisfies the constraints of test cases $15, 16$.

\subsection*{Constraints}

For all test data, it is guaranteed that:
\begin{itemize}
\item $1 \le t \le 5$;
\item $2 \le n \le 10^5$, and $n$ is even;
\item for all $1 \le i \le n$, $1 \le j \le 3$, we have
  $0 \le a_{i,j} \le 2 \times 10^4$.
\end{itemize}

\begin{center}
\begin{tabular}{|c|c|c|}
\hline
Test case no. & $n =$ & Special property \\
\hline
$1$ & $2$ & \multirow{4}{*}{None} \\
\cline{1-2}
$2$ & $4$ & \\
\cline{1-2}
$3, 4$ & $10$ & \\
\cline{1-2}
$5 \sim 8$ & $30$ & \\
\hline
$9$ & \multirow{2}{*}{$200$} & B \\
\cline{1-1}\cline{3-3}
$10, 11$ & & None \\
\hline
$12$ & \multirow{4}{*}{$10^5$} & A \\
\cline{1-1}\cline{3-3}
$13, 14$ & & B \\
\cline{1-1}\cline{3-3}
$15, 16$ & & C \\
\cline{1-1}\cline{3-3}
$17 \sim 20$ & & None \\
\hline
\end{tabular}
\end{center}

Special property A: for all $1 \le i \le n$, we have $a_{i,2} = a_{i,3} = 0$.

Special property B: for all $1 \le i \le n$, we have $a_{i,3} = 0$.

Special property C: for all $1 \le i \le n$, $1 \le j \le 3$, each $a_{i,j}$
is generated \textbf{independently and uniformly at random} in
$[0, 2 \times 10^4]$.

\clearpage
\section*{Problem 2: Road Repair (\cn{道路修复})
  \hfill{\normalsize\href{https://qoj.ac/problem/14598}{QOJ 14598}}}

\noindent
Time limit: 1 s. \quad Memory limit: 512 MB. \quad Total points: 100.

\subsection*{Problem description}

The transportation system of Country C consists of $n$ cities and $m$
bidirectional roads, each connecting two cities.
The $i$-th ($1 \le i \le m$) road connects cities $u_i$ and $v_i$.
\textbf{Any two cities can reach each other through some of the roads.}

However, recently, because of a major earthquake, all $m$ roads were
destroyed, and the cost of repairing the $i$-th ($1 \le i \le m$) road is
$w_i$.
At the same time, Country C also has $k$ towns that are prepared to undergo
urbanization.
For the $j$-th ($1 \le j \le k$) town, the cost for Country C to urbanize it
is $c_j$.
After the $j$-th ($1 \le j \le k$) town has been urbanized, several roads may
be built between this town and the original $n$ cities, where the cost of
building a road between it and the $i$-th ($1 \le i \le n$) city is
$a_{j,i}$.
Country C may choose \textbf{any number} of these $k$ towns to urbanize, and
may also choose not to urbanize any town.

To restore the traffic between the cities as soon as possible, the
government of Country C wants to connect the \textbf{original} $n$ cities
pairwise at the lowest cost, that is, any two of the original cities can
reach each other through some of the repaired or newly built roads.
You need to help them find the minimum cost of connecting the original $n$
cities pairwise.

\subsection*{Input}

The first line of the input contains three non-negative integers $n, m, k$,
the number of original cities, the number of roads, and the number of towns
prepared to undergo urbanization, respectively.

Line $i+1$ ($1 \le i \le m$) of the input contains three non-negative
integers $u_i, v_i, w_i$, the two cities connected by the $i$-th road and
the cost of repairing this road.

Line $j+m+1$ ($1 \le j \le k$) of the input contains $n+1$ non-negative
integers $c_j, a_{j,1}, a_{j,2}, \dots, a_{j,n}$, the cost of urbanizing
the $j$-th town and the costs of building roads between this town and the
original cities, respectively.

\subsection*{Output}

Output one line containing one non-negative integer, the minimum cost of
connecting the original $n$ cities pairwise.

\subsection*{Sample 1 input}

\begin{verbatim}
4 4 2
1 4 6
2 3 7
4 2 5
4 3 4
1 1 8 2 4
100 1 3 2 4
\end{verbatim}

\subsection*{Sample 1 output}

\begin{verbatim}
13
\end{verbatim}

\subsection*{Explanation of sample 1}

The government of Country C can choose to repair the $3$rd and $4$th roads,
then urbanize the $1$st town and build the roads between it and cities
$1, 3$.
The total cost is $5 + 4 + 1 + 1 + 2 = 13$.
It can be shown that there is no cost smaller than $13$ that connects the
original $4$ cities pairwise.

\subsection*{Sample 2}

Large sample, not printed here;
provided as the downloadable files \texttt{road/road2.in} and \texttt{road/road2.ans}.

This sample satisfies the constraints of test cases $11, 12$.

\subsection*{Sample 3}

Large sample, not printed here;
provided as the downloadable files \texttt{road/road3.in} and \texttt{road/road3.ans}.

This sample satisfies the constraints of test cases $13, 14$.

\subsection*{Sample 4}

Large sample, not printed here;
provided as the downloadable files \texttt{road/road4.in} and \texttt{road/road4.ans}.

This sample satisfies the constraints of test cases $15, 16$.

\subsection*{Constraints}

For all test data, it is guaranteed that:
\begin{itemize}
\item $1 \le n \le 10^4$, $1 \le m \le 10^6$, $0 \le k \le 10$;
\item for all $1 \le i \le m$, we have $1 \le u_i, v_i \le n$, $u_i \ne v_i$
  and $0 \le w_i \le 10^9$;
\item for all $1 \le j \le k$, we have $0 \le c_j \le 10^9$;
\item for all $1 \le j \le k$, $1 \le i \le n$, we have
  $0 \le a_{j,i} \le 10^9$;
\item any two of the original cities can reach each other through some of
  the original roads.
\end{itemize}

\begin{center}
\begin{tabular}{|c|c|c|c|c|}
\hline
Test case no. & $n \le$ & $m \le$ & $k \le$ & Special property \\
\hline
$1 \sim 4$ & $10^4$ & $10^6$ & $0$ & None \\
\hline
$5, 6$ & \multirow{6}{*}{$10^3$} & \multirow{2}{*}{$10^5$}
  & \multirow{4}{*}{$5$} & A \\
\cline{1-1}\cline{5-5}
$7, 8$ & & & & None \\
\cline{1-1}\cline{3-3}\cline{5-5}
$9, 10$ & & \multirow{7}{*}{$10^6$} & & A \\
\cline{1-1}\cline{5-5}
$11, 12$ & & & & None \\
\cline{1-1}\cline{4-5}
$13, 14$ & & & \multirow{2}{*}{$10$} & A \\
\cline{1-1}\cline{5-5}
$15, 16$ & & & & None \\
\cline{1-2}\cline{4-5}
$17, 18$ & \multirow{3}{*}{$10^4$} & & \multirow{2}{*}{$5$} & A \\
\cline{1-1}\cline{5-5}
$19, 20$ & & & & \multirow{2}{*}{None} \\
\cline{1-1}\cline{4-4}
$21 \sim 25$ & & & $10$ & \\
\hline
\end{tabular}
\end{center}

Special property A: for all $1 \le j \le k$, we have $c_j = 0$, and there
exists $1 \le i \le n$ such that $a_{j,i} = 0$.

\clearpage
\section*{Problem 3: Homophonic Replacement (\cn{谐音替换})
  \hfill{\normalsize\href{https://qoj.ac/problem/14599}{QOJ 14599}}}

\noindent
Time limit: 1 s. \quad Memory limit: 2048 MB. \quad Total points: 100.

\subsection*{Problem description}

Little W is a competitive programmer who likes linguistics.
In linguistics, homophonic replacement means replacing original words by
words with the same or a similar pronunciation.
Little W found that the process of homophonic replacement can be described
with strings.
Specifically, Little W defines homophonic replacement as the following
string problem.

Given $n$ pairs of strings, the $i$-th ($1 \le i \le n$) pair of strings is
$(s_{i,1}, s_{i,2})$, satisfying $|s_{i,1}| = |s_{i,2}|$, where $|s|$
denotes the length of the string $s$.

For a string $s$, a \textbf{replacement} of $s$ is defined as follows:
\begin{itemize}
\item For some substring $y$ of $s$, if there exists $1 \le i \le n$ such
  that $y = s_{i,1}$, then replace $y$ by $y' = s_{i,2}$.
  Specifically, let $s = x + y + z$, where $x$ and $z$ may be empty and
  ``$+$'' denotes string concatenation; then the replacement of $s$ yields
  the string $s' = x + y' + z$.
\end{itemize}

Little W poses $q$ questions.
The $j$-th ($1 \le j \le m$) question gives two \textbf{distinct} strings
$t_{j,1}, t_{j,2}$, and she wants to know how many replacements of the
string $t_{j,1}$ yield the string $t_{j,2}$.
Two replacements of $s$ are different if and only if
\textbf{the position of the substring $y$ is different or the pair
$(s_{i,1}, s_{i,2})$ used for the replacement is different}, that is,
$x, z$ are different or $i$ is different.
You need to answer all the questions posed by Little W.

\subsection*{Input}

The first line of the input contains two positive integers $n, q$, the
number of pairs of strings and the number of questions posed by Little W,
respectively.

Line $i+1$ ($1 \le i \le n$) of the input contains two strings
$s_{i,1}, s_{i,2}$, the $i$-th pair of strings.

Line $j+n+1$ ($1 \le j \le q$) of the input contains two strings
$t_{j,1}, t_{j,2}$, the $j$-th question posed by Little W.

\subsection*{Output}

Output $q$ lines, where the $j$-th ($1 \le j \le q$) line contains one
non-negative integer, the number of replacements of the string $t_{j,1}$
that yield the string $t_{j,2}$.

\subsection*{Sample 1 input}

\begin{verbatim}
4 2
xabcx xadex
ab cd
bc de
aa bb
xabcx xadex
aaaa bbbb
\end{verbatim}

\subsection*{Sample 1 output}

\begin{verbatim}
2
0
\end{verbatim}

\subsection*{Explanation of sample 1}

For Little W's first query, there are $2$ replacements of $t_{1,1}$ that
yield $t_{1,2}$:
\begin{enumerate}
\item Let $x, z$ both be the empty string, $y = \texttt{xabcx}$, $i = 1$;
  then $y' = \texttt{xadex}$, and the replacement yields \texttt{xadex};
\item Let $x = \texttt{xa}$, $y = \texttt{bc}$, $z = \texttt{x}$, $i = 3$;
  then $y' = \texttt{de}$, and the replacement yields \texttt{xadex}.
\end{enumerate}

\subsection*{Sample 2 input}

\begin{verbatim}
3 4
a b
b c
c d
aa bb
aa b
a c
b a
\end{verbatim}

\subsection*{Sample 2 output}

\begin{verbatim}
0
0
0
0
\end{verbatim}

\subsection*{Sample 3}

Large sample, not printed here;
provided as the downloadable files \texttt{replace/replace3.in} and \texttt{replace/replace3.ans}.

This sample satisfies the constraints of test cases $11, 12$.

\subsection*{Sample 4}

Large sample, not printed here;
provided as the downloadable files \texttt{replace/replace4.in} and \texttt{replace/replace4.ans}.

This sample satisfies the constraints of test cases $15, 16$.

\subsection*{Constraints}

Let $L_1 = \sum_{i=1}^{n} |s_{i,1}| + |s_{i,2}|$ and
$L_2 = \sum_{j=1}^{q} |t_{j,1}| + |t_{j,2}|$.
For all test data, it is guaranteed that:
\begin{itemize}
\item $1 \le n, q \le 2 \times 10^5$;
\item $2 \le L_1, L_2 \le 5 \times 10^6$;
\item for all $1 \le i \le n$, $s_{i,1}, s_{i,2}$ consist only of lowercase
  English letters, and $|s_{i,1}| = |s_{i,2}|$;
\item for all $1 \le j \le q$, $t_{j,1}, t_{j,2}$ consist only of lowercase
  English letters, and $t_{j,1} \ne t_{j,2}$.
\end{itemize}

\begin{center}
\begin{tabular}{|c|c|c|c|}
\hline
Test case no. & $n, q \le$ & $L_1, L_2 \le$ & Special property \\
\hline
$1, 2$ & $10^2$ & $200$ & \multirow{2}{*}{None} \\
\cline{1-3}
$3 \sim 5$ & \multirow{2}{*}{$10^3$} & $2{,}000$ & \\
\cline{1-1}\cline{3-4}
$6$ & & \multirow{3}{*}{$10^6$} & AB \\
\cline{1-2}\cline{4-4}
$7, 8$ & $10^4$ & & A \\
\cline{1-2}\cline{4-4}
$9, 10$ & \multirow{5}{*}{$2 \times 10^5$} & & B \\
\cline{1-1}\cline{3-4}
$11, 12$ & & $2 \times 10^6$ & None \\
\cline{1-1}\cline{3-4}
$13, 14$ & & \multirow{3}{*}{$5 \times 10^6$} & A \\
\cline{1-1}\cline{4-4}
$15, 16$ & & & B \\
\cline{1-1}\cline{4-4}
$17 \sim 20$ & & & None \\
\hline
\end{tabular}
\end{center}

Special property A: $q = 1$.

Special property B: define a string $s$ to be \textbf{special} if and only if
$s$ consists only of the characters \texttt{a} and \texttt{b}, and the
character \texttt{b} occurs \textbf{exactly} once in $s$.
For all $1 \le i \le n$, $s_{i,1}, s_{i,2}$ are special, and for all
$1 \le j \le q$, $t_{j,1}, t_{j,2}$ are special.

\clearpage
\section*{Problem 4: Employee Recruitment (\cn{员工招聘})
  \hfill{\normalsize\href{https://qoj.ac/problem/14600}{QOJ 14600}}}

\noindent
Time limit: 1 s. \quad Memory limit: 512 MB. \quad Total points: 100.

\subsection*{Problem description}

Little Z and Little H want to start a company together.
A total of $n$ people, numbered $1 \sim n$, have come to apply.
Little Z and Little H want to hire at least $m$ people.

Little H is the interviewer and will interview one person per day over the
next $n$ days.
Little Z is in charge of deciding the order in which the applicants come
for their interviews.
Specifically, Little Z may choose a permutation $p$ of $1 \sim n$, and then
on day $i$ ($1 \le i \le n$) notify the person numbered $p_i$ to come for
an interview.

Little H has prepared $n$ sets of interview questions of varying difficulty.
Since the $n$ applicants are of roughly the same ability, for the same set
of questions, the results of all people are identical.
Specifically, the difficulty of the interview questions on day $i$
($1 \le i \le n$) is $s_i \in \{0, 1\}$, where $s_i = 0$ means that this
set is hard and nobody can solve it, and $s_i = 1$ means that this set is
easy and everybody can solve it.
Little H decides whether to hire according to the interviewee's result,
that is, if the interviewee does not solve the questions, they are
rejected; otherwise they are hired.

However, everybody's patience has a certain limit; if too many people were
not hired before their own interview, they will simply give up on attending
the interview.
Specifically, the patience limit of the person numbered $i$
($1 \le i \le n$) is described by a non-negative integer $c_i$:
if \textbf{at least} $c_i$ people before them have been rejected or have
given up on attending the interview, then they will also give up on
attending the interview.

Little Z wants to know how many interview orders $p$ in total allow them to
hire at least $m$ people.
You need to help Little Z find the number of permutations $p$ for which at
least $m$ people are hired.
Since the answer may be large, you only need to find the answer modulo
$998{,}244{,}353$.

\subsection*{Input}

The first line of the input contains two positive integers $n, m$, the
number of applicants and the number of people they want to hire,
respectively.

The second line of the input contains a string $s_1 \dots s_n$ of length
$n$, the difficulty of the interview questions on each day.

The third line of the input contains $n$ non-negative integers
$c_1, c_2, \dots, c_n$, the patience limit of each person.

\subsection*{Output}

Output one line containing one non-negative integer, the number of
permutations $p$ for which at least $m$ people are hired, modulo
$998{,}244{,}353$.

\subsection*{Sample 1 input}

\begin{verbatim}
3 2
101
1 1 2
\end{verbatim}

\subsection*{Sample 1 output}

\begin{verbatim}
2
\end{verbatim}

\subsection*{Explanation of sample 1}

There are the following $2$ interview orders $p$ that allow Little Z and
Little H to hire at least $2$ people:
\begin{enumerate}
\item $p = [1, 2, 3]$, hiring the person numbered $1$ and then the person
  numbered $3$;
\item $p = [2, 1, 3]$, hiring the person numbered $2$ and then the person
  numbered $3$.
\end{enumerate}

\subsection*{Sample 2 input}

\begin{verbatim}
10 5
1101111011
6 0 4 2 1 2 5 4 3 3
\end{verbatim}

\subsection*{Sample 2 output}

\begin{verbatim}
2204128
\end{verbatim}

\subsection*{Sample 3}

Large sample, not printed here;
provided as the downloadable files \texttt{employ/employ3.in} and \texttt{employ/employ3.ans}.

This sample satisfies the constraints of test cases $6 \sim 8$.

\subsection*{Sample 4}

Large sample, not printed here;
provided as the downloadable files \texttt{employ/employ4.in} and \texttt{employ/employ4.ans}.

This sample satisfies the constraints of test cases $12 \sim 14$.

\subsection*{Sample 5}

Large sample, not printed here;
provided as the downloadable files \texttt{employ/employ5.in} and \texttt{employ/employ5.ans}.

This sample satisfies the constraints of test cases $18 \sim 21$.

\subsection*{Constraints}

For all test data, it is guaranteed that:
\begin{itemize}
\item $1 \le m \le n \le 500$;
\item for all $1 \le i \le n$, we have $s_i \in \{0, 1\}$;
\item for all $1 \le i \le n$, we have $0 \le c_i \le n$.
\end{itemize}

\begin{center}
\begin{tabular}{|c|c|c|c|}
\hline
Test case no. & $n \le$ & $m$ & Special property \\
\hline
$1, 2$ & $10$ & \multirow{4}{*}{$\le n$} & \multirow{2}{*}{None} \\
\cline{1-2}
$3 \sim 5$ & $18$ & & \\
\cline{1-2}\cline{4-4}
$6 \sim 8$ & \multirow{2}{*}{$10^2$} & & \multirow{2}{*}{A} \\
\cline{1-1}
$9 \sim 11$ & & & \\
\hline
$12 \sim 14$ & \multirow{5}{*}{$500$} & $= 1$ & \multirow{2}{*}{None} \\
\cline{1-1}\cline{3-3}
$15$ & & $= n$ & \\
\cline{1-1}\cline{3-4}
$16, 17$ & & \multirow{3}{*}{$\le n$} & A \\
\cline{1-1}\cline{4-4}
$18 \sim 21$ & & & B \\
\cline{1-1}\cline{4-4}
$22 \sim 25$ & & & None \\
\hline
\end{tabular}
\end{center}

Special property A: for all $1 \le i \le n$, we have $s_i = 1$.

Special property B: among $s_1, s_2, \dots, s_n$ at most $18$ take the
value $1$, that is, $\sum_{i=1}^{n} s_i \le 18$.

\section*{Source}

CSP-S 2025 (CCF CSP-J/S 2025 second round, Senior level), held 2025-10-25.
The statements were taken from the official statement PDFs embedded on the
QOJ problem pages.
Contest page: \texttt{https://qoj.ac/contest/2571}.
Problem pages:
\begin{itemize}
\item Problem 1: \texttt{https://qoj.ac/problem/14597}
\item Problem 2: \texttt{https://qoj.ac/problem/14598}
\item Problem 3: \texttt{https://qoj.ac/problem/14599}
\item Problem 4: \texttt{https://qoj.ac/problem/14600}
\end{itemize}

\end{document}
